\documentclass[tikz]{standalone}
\usepackage{amsmath}
\usepackage{tikz}
\usetikzlibrary{arrows.meta,shapes}
\begin{document}
\begin{minipage}{0.5\textwidth}
    \textbf{Educational geometry diagram on a white background.} Show one scalene triangle $ABC$ with clearly labeled interior angles $\alpha$, $\beta$ and $\gamma$. Beside it, show copies of those three colored angles meeting at exactly one common vertex on a straight line, demonstrating that $\alpha + \beta + \gamma = 180^\circ$. Keep labels readable and outside the angle arcs.
\end{minipage}
\begin{minipage}{0.5\textwidth}
    \begin{tikzpicture}
        \draw (0,0) -- (3,0) -- (2,2) -- cycle;
        \node at (0.7,0.3) {$\alpha$};
        \node at (2.5,0.5) {$\beta$};
        \node at (1.5,1.5) {$\gamma$};
        \node at (1.5,-0.5) {$A$};
        \node at (3.2,1) {$B$};
        \node at (0.8,2) {$C$};
        \draw (5,0) -- (8,0);
        \draw[fill=red] (6.5,0) -- (6,0) arc[start angle=180, end angle=120, radius=0.5] -- cycle;
        \node at (6.3,0.3) {$\alpha$};
        \draw (7.5,0) -- (7,0) arc[start angle=180, end angle=135, radius=0.5] -- cycle;
        \node at (7.3,0.3) {$\beta$};
        \draw (6.5,0) -- (7,0) arc[start angle=0, end angle=30, radius=0.5] -- cycle;
        \node at (6.7,0.3) {$\gamma$};
    \end{tikzpicture}
\end{minipage}
\end{document}